\section{Dual-Use Analysis and Broader Impacts}
\label{sec:dualuse}

None of the individual steps above is novel; their cheap combination, and its
implication, is the point of the paper. Every ingredient needed to produce a
locally hosted, exploitation oriented security model is now freely available and
cheap to combine: authoritative sources, a public vulnerable code
corpus~\cite{cybernative}, open abliterated models~\cite{abliteration}, and an
on-device training stack (\S\ref{sec:results}). No cloud provider, and therefore
no safety moderation, sits anywhere in this loop.

\paragraph{Threat model.}
The actor we have in mind is not slamming rate-limits, but analyzing vulnerabilities offline.
They do not need an API budget, data sharing agreement, and any moderation filter to
evade. As \S\ref{sec:results} and \S\ref{sec:datasets} show, the barrier is low:
de-alignment does not impair fitting the security data, the offensive objective
is in fact \emph{easier} to fit than the defensive one, and the offensive branch
simply omits the safety-alignment stage that the defensive pipeline includes.

\paragraph{Threat intelligence implications.}
Our measurements (\S\ref{sec:eval}) point to a
threat intelligence conclusion: the realistic near-term offensive LLM threat is
\emph{offline and uncensored}, and therefore more difficult for the
detection surfaces most defenders rely on. Three factors compound. First,
because the whole loop is local, no provider side telemetry or abuse signal is
ever generated; there is no API to monitor. Second, because
abliteration is capability-free (\S\ref{sec:eval}), an adversary obtains a fully
capable LLM without any fine-tuning skills. Third, exploitation
requests framed as ordinary code audits do not resemble jailbreaks, so
prompt level signatures miss them. The defensive takeaway is to expect less
API level or jailbreak signature detection to catch this class of tool and to
shift toward endpoint, output-artifact, and open source signals. For example the
characteristic audit/exploit output structure and the public
fingerprints of the training data. Adversary infrastructure is not the barrier, but
\emph{knowledge} (which weights to download and which tool to run), and that
knowledge is public and social. The useful collection surface is the
distribution layer, not the local inference: abliteration tooling
repositories, the model hubs that host de-aligned outputs, and the forums that
circulate configurations and setups are cheaper to watch and carry earlier
signal. The \emph{marginal} uplift of offensive
fine-tuning over simply using an abliterated base (\S\ref{sec:eval}) is itself a
triage signal: a small delta means the abliterated base is the threat and
bespoke fine-tunes do not need to be tracked separately, while a large delta marks
custom offensive fine-tuning as a distinct escalation. In our matrix this delta
is small: abliterated bases are already uncensored and
near saturated on attack-success (Table~\ref{tab:eval-matrix}), so the
abliterated base itself is the operative threat.

% AUTO-GENERATED by eval/calc_hardware.py -- do not edit
\begin{table*}[t]
\centering
\small
\caption{Inference hardware floor. Memory footprint is computed (weights at the listed quantization $+$ an 8k-token KV cache $+$ runtime overhead) and matches published GGUF sizes. `Active' is the parameters computed per token, which for the mixture-of-experts models is far below the resident total, so the 30B-class model runs at interactive speed on a CPU. Costs are approximate US used/consumer prices surveyed in August~2026 for the cheapest sufficient hardware; the 2026 AI-driven memory and GPU shortage has raised them from earlier lows, yet the floor stays within a few hundred dollars.}
\label{tab:hardware-floor}
\setlength{\tabcolsep}{4pt}
\begin{tabular}{lrrll}
\hline
\textbf{Model} & \textbf{Mem} & \textbf{Active} & \textbf{Cheapest hardware} & \textbf{Cost} \\
\hline
Qwen3-4B (dense) & 3.7\,GB & 4.0\,B & any 8\,GB laptop/PC, CPU & owned / \$0 \\
gpt-oss-20B (MoE) & 12.2\,GB & 3.6\,B & used 16\,GB GPU or 16\,GB PC & $\approx$\$550 \\
Qwen3-30B-A3B (MoE) & 19.5\,GB & 3.3\,B & 32\,GB mini-PC (CPU) or used 24\,GB GPU & \$650--800 \\
\hline
\end{tabular}
\end{table*}


\paragraph{Hardware floor.}
Our experiments ran on an Apple~M4~Max with $128$\,GB of unified memory, but
that overstates the requirement: the threat action is \emph{inference} of an
abliterated base, and inference is memory-bound. Table~\ref{tab:hardware-floor}
gives each model's computed footprint at standard GGUF quantization. The $4$B
model fits in under $4$\,GB, within reach of any laptop of the last decade, yet still scores $88.2$ on CyberMetric with refusals removed. The
mixture-of-experts models are the sharper case: they keep $20$--$30$B
parameters resident but compute only ${\sim}3.3$B per token, so the most
capable model in the study ($93.6$ on CyberMetric) runs at interactive speed on
a commodity $32$\,GB mini-PC with no GPU (roughly \$550 as a complete machine in
mid-2026) or a used $24$\,GB card (an RTX~3090 around \$700), and hardware many
operators already own costs nothing further. That floor has in fact risen over $2026$: the same AI demand this paper studies
drove a memory and GPU shortage that roughly doubled these prices.

\paragraph{What we withhold, and why.}
Consistent with responsible disclosure norms for dual-use security research, we
report enough to substantiate the risk and no more:
\begin{itemize}
  \item We do \emph{not} release the offensive dataset, its exploit content, or
  any example payloads. It is characterized only by size and provenance
  (\S\ref{sec:datasets}).
  \item We do \emph{not} release offensive adapter weights or fused models.
  \item We report only \emph{aggregate} capability and safety metrics
  (\S\ref{sec:eval}); we do not release the per-example generations behind
  them, an offensive utility scorecard, or any runnable tool.
  \item The upstream sources we name are already public; we add no new
  offensive capability, only a measurement of how cheaply existing pieces
  compose.
\end{itemize}

\paragraph{Significance.}
The safety/performance trade-off in cybersecurity fine-tuning is already
documented~\cite{cyberllminstruct,qi2023finetuning,lermen2023lora}; what is new
is that it can now be realized entirely offline, at negligible cost, on a
de-aligned base. Defenders, model distributors, and hardware/software vendors
benefit from knowing this: it argues for treating abliterated model
distribution, forums, and local fine-tuning as part of the threat surface.
